\chapter{Introduction}
\label{chap:introduction2}

Neutron stars are compact stellar remnants formed in core–collapse supernovae,
with typical masses of order one to two solar masses and radii of roughly ten
kilometres. Observations accumulated over the past decades—most notably precise
mass measurements from radio timing of pulsars, X-ray observations of neutron-star
surfaces, and gravitational-wave signals from neutron-star mergers—have established
neutron stars as key laboratories for physics at extreme densities. Any theoretical
description of these objects must therefore be consistent with a growing set of
astrophysical constraints \cite{LattimerPrakash2004,glendenning2000}.

The structure of a non-rotating neutron star is described within general relativity
by the Tolman–Oppenheimer–Volkoff (TOV) equations \cite{tolman1939,oppenheimer1939}.
These equations relate the pressure, energy density, and enclosed mass of the star
under the assumption of hydrostatic equilibrium. In order to solve them, one must
specify an equation of state (EoS) for dense matter. The EoS encodes assumptions
about the microscopic constituents of the star and their interactions, and it is
the primary source of theoretical uncertainty in neutron-star modelling.

Early neutron-star models employed simple equations of state based on ideal,
degenerate Fermi gases of neutrons. These models provide a useful starting point
and lead to the classic Oppenheimer–Volkoff limit of approximately
$0.7\,M_{\odot}$ for the maximum neutron-star mass \cite{oppenheimer1939}. While
historically important, this result is now known to be far below the masses of
observed neutron stars, demonstrating that interactions between particles are
essential for a realistic description of dense matter.

Including interactions between nucleons leads to significantly stiffer equations
of state and substantially higher maximum masses. Models that treat neutrons and
protons as interacting particles, supplemented by leptons to ensure charge
neutrality and equilibrium, represent a clear improvement over idealized Fermi-gas
descriptions. Such interacting models are capable of producing neutron stars with
masses closer to those observed in nature. However, even these more realistic
descriptions remain uncertain at the highest densities and often struggle to
reproduce the full range of current observational constraints. This issue was
demonstrated explicitly in the project thesis (Part~I), where we looked at the
$npe\mu$ model for matter. It was shown that,
despite significant improvement over idealized models, the predicted mass-radius
still do not align with well-measured neutron stars.

The persistence of these discrepancies suggests that further improvements to the
equation of state are required. One possible direction is that the relevant
degrees of freedom at the highest densities differ qualitatively from those of
ordinary nuclear matter. The development of quantum chromodynamics (QCD) in the
early 1970s, as the fundamental theory of the strong interaction, provided a
microscopic description of hadrons in terms of quarks and gluons
\cite{GrossWilczek1973}. This, in turn, opened the possibility of considering
regimes in which quark degrees of freedom are no longer confined inside hadrons,
but appear explicitly in bulk matter.

Shortly after the formulation of QCD emerged the idea that
matter at sufficiently high density might undergo a transition from a hadronic
description to one in terms of deconfined quarks. Early studies in the 1970s and
1980s examined the implications of such a transition for compact stars, leading
to the proposal of new classes of stellar objects
\cite{Itoh1970,Bodmer1971,Witten1984}. Within this framework, neutron stars may
contain cores of deconfined quark matter, commonly referred to as hybrid stars,
or, in more speculative scenarios, be composed entirely of quark matter, giving
rise to quark or strange stars.

A theoretical description of these possibilities requires a framework that goes
beyond phenomenological models of hadronic matter. Quantum field theory provides
the natural language for matter at a fundamental
level, with QCD as the underlying theory governing the
interactions between quarks and gluons \cite{PeskinSchroeder1995}. While QCD is
well established as the correct microscopic theory of the strong interaction,
its direct application to the high-density, low-temperature regime relevant for
neutron stars remains notoriously challenging. As a result, first-principles
calculations in this domain are limited, and one must instead rely on
approximate approaches. This has motivated the development of effective theories
inspired by QCD, which aim to capture the relevant degrees of freedom and
symmetry properties of strongly interacting matter while remaining tractable
\cite{glendenning2000}. Such approaches provide a practical starting point for
studying deconfined quark matter in compact stars and for exploring how a change
in the microscopic description of matter can affect macroscopic stellar
properties.

The present master thesis follows this line of reasoning. Motivated by the
limitations of hadronic equations of state and supported by the findings of
Part~I, it explores quantum field–theoretic descriptions of dense matter with
quark degrees of freedom and investigates their implications for the structure
and observable properties of compact stars.
